\section{Introduction}
\label{sec:intro}

Benchmarks shape what the field optimizes. As frontier models saturate graduate-level multiple-choice suites (GPT-5.2 reports 92\% on GPQA~\cite{rein2023gpqa}; see also OpenAI's FrontierScience motivation~\cite{openai2025frontierscience}), the community increasingly turns to domain-specific, quantitatively verifiable evaluations---SWE-bench for software engineering~\cite{jimenez2023swebench}, SciBench and PHYSICS for college and graduate physics~\cite{wang2023scibench,feng2025physics}, and rubric-graded expert tasks in FrontierScience. Nuclear fusion is a conspicuous gap. Fusion physics---especially inertial confinement fusion (ICF), whose 2022 ignition milestone at the National Ignition Facility (NIF) made front-page news worldwide~\cite{llnl2022ignition}---combines radiation hydrodynamics, plasma kinetics, laser--plasma instabilities, and high-energy-density equations of state. It is a natural stress test for physical AI: problems have exact or semi-analytic answers, unit systems (cgs vs.\ SI) trip models constantly, and multi-step energy-accounting chains punish shallow pattern matching.

Yet the existing landscape offers almost nothing quantitative for fusion. General science benchmarks include essentially zero fusion problems (Section~\ref{sec:related}). The only fusion-specific benchmark we are aware of, FusionBench (PKU-XLab)~\cite{pkufusionbench}, contains 10{,}605 knowledge questions (single/multiple choice and true/false); our audit of its public dataset snapshot finds only 80 problems (0.75\%) containing any numeric value with a physical unit, and no independent answer-verification process. Competition-style fusion efforts---the NeurIPS 2026 Fusion Equilibrium Challenge~\cite{sophelio2026fec} and the TokaMark benchmarks~\cite{gupta2026tokamark,ibm2026tokamark}---target machine-learning models for tokamak diagnostics and equilibrium reconstruction, not physics reasoning. To our knowledge, no benchmark evaluates quantitative fusion-physics reasoning---numeric computation, formula derivation, code---with an ICF focus.

\textbf{Contributions.}
\begin{enumerate}
    \item \textbf{Physical-AI-Fusion} (v0.1): 106 verified problems across 13 topics with ICF at 52\% (ignition \& gain, Rayleigh--Taylor/Richtmyer--Meshkov instabilities, laser--plasma interactions, radiation hydrodynamics, EOS \& opacity, implosion dynamics, target design, diagnostics), plus magnetic confinement, Z-pinch, plasma fundamentals, fusion engineering, and computational physics. Task types: 62 numeric, 30 concept, 13 derivation (rubric-graded), 1 figure.
    \item \textbf{A mechanical verification pipeline} (Section~\ref{sec:verification}): schema validation, sandboxed independent recomputation of every numeric answer, dual verification methods for non-numeric answers, and a staging$\rightarrow$verified promotion gate with a public log. The pipeline has already intercepted two constructor arithmetic slips and rejected one local dataset wholesale---evidence it functions as a correctness gate rather than a rubber stamp.
    \item \textbf{A SWE-bench-style harness} (Section~\ref{sec:harness}): \texttt{pip install} plus a single \texttt{evaluate} command yields deterministic scores with per-topic/difficulty breakdowns; a verified split is the headline set; a Dockerfile reproduces the environment.
    \item \textbf{Baselines and analysis} (Section~\ref{sec:experiments}): deterministic grading on the verified split, with failure-mode analysis by topic and difficulty.
\end{enumerate}

We release the dataset (CC BY 4.0), harness and scripts (MIT), verification log, and this paper's source at \url{https://github.com/NiubilityZXC/physical-ai-fusion}.
